% !TeX root = BookN.tex

\title{Stress State of Elastic Bodies With Cracks Under Harmonic Oscillations}
\titlerunning{Stress State of Elastic Bodies With Cracks}

\author{Vsevolod Popov and Olga Kyrylova}

\institute{%
	Vsevolod Popov
	\at  \AffOdesaMaritime
	\\
	\email{dr.vg.popov@gmail.com}
	\and
	Olga Kyrylova
	\at \AffOdesaMaritime
	\\
	\email{olga.i.kyrylova@gmail.com}
}

\graphicspath{{31_Popov/}}

\maketitle


\abstract{
	This paper presents numerical-analytical methods for solving problems related to stress determination, the calculation of stress intensity factors (SIFs), and the investigation of resonance phenomena in elastic bodies containing systems of internal tunnel cracks. Both unbounded elastic media and cylindrical bodies of arbitrary cross-section are considered. Under conditions of longitudinal shear and plane strain, the corresponding fracture mechanics problems are reduced to two-dimensional equations of motion for an elastic medium in plane domains bounded by arbitrary smooth curves and containing straight-line cracks.
	The proposed solution approach is based on representing the displacement fields using discontinuous solutions of the governing equations that incorporate displacement jumps along the crack surfaces. In the case of an unbounded medium, the total displacement is constructed as a superposition of such discontinuous solutions, with the unknown jumps determined from integro-differential equations derived from boundary conditions on the cracks. These equations are solved efficiently using the method of mechanical quadratures over a broad frequency range. For systems with many cracks, an iterative method is developed that avoids the need to solve large-scale integro-differential systems. Instead, at each iteration, a set of independent equations is solved for each individual crack.
		For cylindrical bodies, the displacement field is represented as the sum of two components: one is the aforementioned superposition of discontinuous solutions, while the other is an unknown function that satisfies the governing equations and enforces boundary conditions on the outer surface of the body. This unknown function is expressed as a linear combination of linearly independent particular solutions with undetermined coefficients. This formulation enables the separate satisfaction of boundary conditions on both the cracks and the external boundary. The integro-differential equations on the cracks are solved numerically using the mechanical quadrature method, and the boundary conditions on the outer surface are enforced by determining the coefficients via the collocation method.
		Numerical examples demonstrate that the mutual interaction of cracks can significantly amplify the stress intensity compared to a single isolated crack. This is manifested in the appearance of absolute maxima in the SIFs, with values several times greater than those observed for a single crack. Additionally, for cylindrical bodies, the method facilitates the analysis of resonance phenomena, which are particularly important in fracture mechanics due to the sharp (formally infinite in linear theory) increase in SIF magnitude at resonance. It is shown that the first resonance frequency is highly sensitive to the location, proximity to boundaries, orientation, and size of the cracks.
}


%\textbf{Key words }Crack, cylindrical bodies, iteration method, dynamic stress intensive factor, resonance phenomena, elastic waves


\section{Introduction}

The relevance of solving dynamic elasticity problems for bodies with cracks lies in the wave reflection phenomena caused by dynamic loading. Waves reflected from crack surfaces significantly alter the local stress distribution and can result in dynamic SIFs that substantially exceed their static counterparts. This enhancement is observed even for a single crack~\cite{Popov13, Popov21, Popov27, Popov28, Popov31}.

The problem of determining the displacement and stress fields generated by waves interacting with a single crack in an elastic medium is currently well understood. However, the situation becomes more complex when multiple interacting cracks are present in the body. Advances in computational power and specialized software tools—such as ANSYS, ABAQUS, NASGROW, and AFCROSS—have enabled the use of direct numerical methods to solve both dynamic and static fracture mechanics problems~\cite{Popov2, Popov7, Popov8, Popov26}. These methods are highly versatile, accommodating bodies and cracks of arbitrary shape without the need for simplifying assumptions.

Nevertheless, as noted in the cited works, the use of direct numerical methods requires substantial mesh refinement, especially near crack tips where stress singularities occur. For systems of closely spaced cracks, mesh refinement must be performed individually for each crack and extended over a large surrounding area, which poses significant computational challenges. These issues limit the practical applicability of finite element and finite difference methods to configurations involving only one or two cracks.

A widely used and effective alternative is the method of boundary integral equations (BIE) and its discrete counterpart, the boundary element method (BEM). The development and application of these methods to dynamic fracture mechanics problems have been extensively studied in monographs and research articles~\cite{Popov28, Popov31, Popov3, Popov5, Popov6, Popov9, Popov11, Popov12, Popov22, Popov25, Popov29}. However, applying BIE methods to bodies with multiple cracks results in large systems of singular integro-differential equations, with the system size scaling with the number of cracks.

To address this challenge, early research focused on dynamically loaded parallel cracks~\cite{Popov10, Popov18, Popov19, Popov23}. Further simplifications and efficiency gains were achieved for periodic arrays of cracks~\cite{Popov4, Popov20, Popov30}, which reduce computational complexity. For general two-dimensional problems involving arbitrarily placed cracks under dynamic loading, several studies~\cite{Popov13, Popov3, Popov9, Popov11, Popov12, Popov29} have reformulated the original boundary-value problems as systems of integral or integro-differential equations. Yet, even in these general frameworks, numerical results are typically limited to configurations with only a few (usually two) cracks.

This motivates the development of a method capable of evaluating dynamic SIFs for systems containing a large number of cracks without solving large-scale systems of integral equations. In the present work, an iterative method is proposed for calculating SIFs in arbitrarily arranged dynamically loaded crack systems. The method involves solving a set of equations for individual cracks at each iteration step, thereby significantly reducing computational complexity.

A characteristic feature of all BIE modifications is the need to enforce boundary conditions both on crack surfaces and on the outer boundary of the body. For finite domains with multiple cracks, this requirement greatly complicates the practical use of BIE techniques. The present study introduces a new method that enables separate and successive satisfaction of boundary conditions on both the cracks and the outer boundary. This approach is employed to determine the stress state of an infinitely long cylinder with cracks and to analyze resonance phenomena arising in such configurations.
eveloped. This method has been used to solve the problems of determining the stressed state of an infinitely long cylinder and analyze the resonance phenomena in the presence of cracks.


\section{Determination of the Stress State and Stress Intensity Factors Under the Action of Longitudinal Shear Waves on a System of Cracks}


Consider an elastic body subjected to longitudinal shear deformation that contains \( N \) thin through-cracks. In the \( Oxy \) plane, these cracks are modeled as line segments of length \( 2d_k \), centered at the points \( (a_k, b_k) \) for \( k = 1, 2, \dots, N \) (Fig.~\ref{PopovFig01}).

\begin{figure}[t]
	\sidecaption[t]
	\includegraphics[width=.64\textwidth]{31_Popov/PopovFig01.jpeg}
	\caption{The system of cracks in the body.}
	\label{PopovFig01}
\end{figure}

A longitudinal shear wave propagates through the body and interacts with the cracks, inducing displacements along the \( Oz \) axis given by
\begin{equation} \label{31.1)} 
	W_{q}(x,y) = A\, e^{i\kappa_2 (x\cos \theta_0 + y\sin \theta_0)}, \quad \kappa_2^2 = \dfrac{\rho \omega^2}{G}, 
\end{equation}
where \( \omega \) is the angular frequency of oscillations, \( \rho \) is the density, and \( G \) is the shear modulus of the material. The angle \( \theta_0 \) denotes the direction of wave propagation with respect to the positive \( Ox \)-axis. The time-dependent factor \( e^{-\omega t} \) is omitted in the following analysis.

Let \( W(x,y) \) denote the \( Oz \)-component of the displacement vector field resulting from waves scattered by the cracks. This displacement must satisfy the Helmholtz equation
\begin{equation} \label{PopovEq625098} 
	\Delta W + \kappa_2^2 W = 0,
\end{equation}
where \( \Delta \) is the Laplace operator in the \( Oxy \) coordinate system.

To formulate boundary conditions on the cracks, we introduce local coordinate systems \( O_k x_k y_k \) associated with each crack. The relationships between the global coordinate system and the local systems, as well as between local systems of different cracks, are given by
\begin{equation} \label{PopovEq491412} 
	\begin{aligned}
		x &= a_k + x_k \cos \alpha_k - y_k \sin \alpha_k, 
		\\
		y &= b_k + x_k \sin \alpha_k + y_k \cos \alpha_k;
		\\
		x_l &= 
		\begin{multlined}[t]
			(a_k - a_l) \cos \alpha_l + (b_k - b_l) \sin \alpha_l 
			\\
			+ x_k \cos(\alpha_k - \alpha_l) - y_k \sin(\alpha_k - \alpha_l);
		\end{multlined}
		\\
		y_l &=
		\begin{multlined}[t]
			-(a_k - a_l)
			+ x_k \sin(\alpha_k - \alpha_l) + y_k \cos(\alpha_k - \alpha_l) 
		\end{multlined}
	\end{aligned} 
\end{equation}
($l,k = 1,2,3,\dots,N$).


Let \( W_{qk}(x_k, y_k) \) and \( W_k(x_k, y_k) \) denote the displacements obtained from \( W_q(x, y) \) and \( W(x, y) \), respectively, after transforming to the local coordinates using formulas \eqref{PopovEq491412}. If the cracks' surfaces are assumed to be traction-free, the following boundary conditions must be satisfied:
\begin{equation} \label{PopovEq757980} 
	\tau_{z y_k}(x_k, 0) = -\tau_{q z y_k}(x_k, 0), \quad -d_k < x_k < d_k \quad (k = 1, 2, \dots, N),
\end{equation}
where
\[
\begin{array}{l} 
	\tau_{z y_k} = G \dfrac{\partial W_k}{\partial y_k}, \quad 
	\tau_{q z y_k} = G \dfrac{\partial W_{qk}}{\partial y_k} = -i \kappa_2 G A \sin(\alpha_k - \theta_0)\, e^{i \kappa_2 Z_{0k}}, 
	\\[0.8em]
	Z_{0k} = a_k \cos \theta_0 + b_k \sin \theta_0 + x_k \cos(\alpha_k - \theta_0) - y_k \sin(\alpha_k - \theta_0).
\end{array}
\]

In addition, the displacements \( W_k(x_k, y_k) \) on the crack surfaces exhibit discontinuities. The corresponding unknown jumps are denoted by
\begin{equation} \label{PopovEq759458} 
	\begin{array}{l}
		W_k(x_k, +0) - W_k(x_k, -0) = \chi_k(x_k),
		\\[0.5em]
		\chi_k(\pm d_k) = 0 \quad (k = 1, 2, \dots, N).
	\end{array}
\end{equation}

With these boundary conditions, the problem is formulated as the determination of the displacement and stress fields of the diffraction wave in the body, as well as the SIFs at the crack tips. To solve this problem, for each crack labeled by index \( l \), discontinuous solutions of equation \eqref{PopovEq625098} are constructed in the corresponding local coordinate system \( O_l x_l y_l \). These solutions incorporate displacement jumps \eqref{PopovEq759458}, which serve as unknown functions to be determined from the boundary conditions \cite{Popov15, Popov17}:
\begin{equation} \label{PopovEq185484} 
	\begin{aligned} 
		W_{d l}(x_l, y_l) &= \dfrac{\partial}{\partial y_l} \int_{-d_l}^{d_l} \chi_l(\eta)\, r_2(\eta - x_l, y_l)\, d\eta \quad (l = 1, 2, \dots, N), 
		\\[0.8em]
		r_2(\eta - x_l, y_l) &= \dfrac{i}{4} H_0^{(1)}\left(\kappa_0 \sqrt{(\eta - x_l)^2 + y_l^2}\right),
	\end{aligned} 
\end{equation}
where \( H_0^{(1)} \) is the Hankel function of the first kind of order zero.

The corresponding stress components are calculated using the formulas:
\begin{equation} \label{PopovEq953094} 
	\begin{aligned} 
		\tau_{d z y_l}(x_l, y_l) &=
		\begin{multlined}[t]
			G \int_{-d_l}^{d_l} \chi_l'(\eta) \frac{\partial}{\partial \eta} r_2(\eta - x_l, y_l)\, d\eta 
			\\
			- \kappa_2^2 G \int_{-d_l}^{d_l} \chi_l(\eta)\, \tau_2(\eta - x_l, y_l)\, d\eta,
		\end{multlined}
		\\
		\tau_{d z x_l}(x_l, y_l) &= 
		G \int_{-d_l}^{d_l} \chi_l'(\eta) \frac{\partial}{\partial y_l} r_2(\eta - x_l, y_l)\, d\eta.
	\end{aligned} 
\end{equation}
The displacement field of the scattered wave in the global coordinate system is then represented as
\begin{equation} \label{PopovEq184833} 
	W(x, y) = \sum_{l = 1}^{N} W_{gl}(x, y),
\end{equation}
where \( W_{gl}(x, y) \) are the contributions obtained from the expressions in \eqref{PopovEq185484} after transforming the coordinates using formulas \eqref{PopovEq491412}.

To use representation \eqref{PopovEq184833}, the unknown displacement jumps defined in \eqref{PopovEq759458} must first be determined. This is achieved by applying the boundary conditions \eqref{PopovEq757980}. For this purpose, the stresses appearing in \eqref{PopovEq757980} are expressed through the stress components in \eqref{PopovEq953094}, followed by a coordinate transformation based on the second group of relations in \eqref{PopovEq491412}. As a result, the following expression is obtained:
\begin{equation} \label{PopovEq529534} 
	\tau_{z y_k} = \sum_{l = 1}^{N} \tau_{kl}, \quad 
	\tau_{kl} = -\tau_{d z x_l} \sin(\alpha_k - \alpha_l) + \tau_{d z y_l} \cos(\alpha_k - \alpha_l).
\end{equation}

\noindent Each term \( \tau_{kl} \) in \eqref{PopovEq529534} is computed using the following expression:
\begin{multline} \label{PopovEq548673} 
	\tau_{kl} = G \int_{-d_l}^{d_l} \chi_l'(\eta) \frac{i \kappa_2 H_1^{(1)}(\kappa_2 R_{kl})}{4 R_{kl}} 
	\\
	\times \left[\eta \cos(\alpha_k - \alpha_l) - x_k - (a_k - a_l) \cos \alpha_k - (b_k - b_l) \sin \alpha_k \right] d\eta 
	\\
	+ \kappa_2^2 G \cos(\alpha_k - \alpha_l) \int_{-d_l}^{d_l} \chi_l(\eta) \frac{i}{4} H_0^{(1)}(\kappa_2 R_{kl})\, d\eta,
\end{multline}
where
\begin{multline*}
	R_{kl}^2 = \Big[\eta - (a_k - a_l) \cos \alpha_l - (b_k - b_l) \sin \alpha_l - x_k \cos(\alpha_k - \alpha_l) 
	\\
	+ y_k \sin(\alpha_k - \alpha_l)\Big]^2 
	+ \Big[-(a_k - a_l) \sin \alpha_l + (b_k - b_l) \cos \alpha_l 
	\\
	+ x_k \sin(\alpha_k - \alpha_l) + y_k \cos(\alpha_k - \alpha_l)\Big]^2.
\end{multline*}

By substituting expressions \eqref{PopovEq529534} and \eqref{PopovEq548673} into the boundary condition \eqref{PopovEq757980}, we obtain a system of integro-differential equations. After isolating the singular parts of the kernels, this system takes the following form:
\begin{multline} \label{PopovEq461086} 
	\frac{1}{2\pi} \int_{-1}^{1} \varphi_k'(\tau) \left[ \frac{1}{\tau - \zeta} + F_{kk}(\tau - \zeta) \right] d\tau 
	\\
	+ \frac{1}{2\pi} \int_{-1}^{1} \varphi_k(\tau) \left[ -\gamma_k^2 \kappa_0^2 \ln |\tau - \zeta| + U_{kk}(\tau - \zeta) \right] d\tau 
	\\
	+ \sum_{\substack{l=1 \\ l \ne k}}^{N} \left[ \frac{1}{2\pi} \int_{-1}^{1} \varphi_l'(\tau) F_{kl}(\tau, \zeta)\, d\tau + \frac{1}{2\pi} \int_{-1}^{1} \varphi_l(\tau) U_{kl}(\tau, \zeta)\, d\tau \right] = f_k(\zeta),
\end{multline}
where \( k = 1, 2, \dots, N \). The kernels of the integral operators in system \eqref{PopovEq461086} are given by the following expressions:

\begin{align*}
	F_{kk}(z) &= O(z \ln |z|), \quad 
	U_{kk}(z) = O(z^2 \ln |z|), \quad z \to 0;
	\\
	F_{kl}(\tau, \zeta) &=
	\begin{multlined}[t]
		\dfrac{i\pi \kappa_0 \gamma_l H_1^{(1)}(\kappa_0 q_{kl})}{q_{kl}} 
		\\
		\times \left( \gamma_l \tau \cos \alpha_{kl} - \gamma_k \zeta - (\varepsilon_k - \varepsilon_l) \cos \alpha_k - (\delta_k - \delta_l) \sin \alpha_k \right),
	\end{multlined}
	\\
	U_{kl}(\tau, \zeta) &= \dfrac{i\pi}{2} \kappa_0^2 \gamma_l^2 \cos \alpha_{kl} H_0^{(1)}(\kappa_0 q_{kl}),
\end{align*}
%
where
\begin{multline*}
	q_{kl}^2 = \gamma_l^2 \tau^2 + \gamma_k^2 \zeta^2 + (\varepsilon_k - \varepsilon_l)^2 + (\delta_k - \delta_l)^2 
	- 2 \gamma_l \gamma_k \tau \zeta \cos \alpha_{kl}
	\\
	- 2 \gamma_l \tau \left[ (\varepsilon_k - \varepsilon_l) \cos \alpha_l + (\delta_k - \delta_l) \cos \alpha_l \right]
	\\
	+ 2 \gamma_k \zeta \left[ (\varepsilon_k - \varepsilon_l) \cos \alpha_k + (\delta_k - \delta_l) \sin \alpha_k \right].
\end{multline*}

In deriving system \eqref{PopovEq461086}, we introduced the following notations:

\begin{align*}
	f_k(\zeta) &= -i \kappa_0 A_0 \sin(\alpha_k - \theta_0)\, e^{i \kappa_0 r_{0k}(\zeta)},
	\\[0.5em]
	r_{0k}(\zeta) &= \varepsilon_k \cos \theta_0 + \delta_k \sin \theta_0 + \gamma_k \zeta \cos(\alpha_k - \theta_0),\quad
	\eta = d_l \tau, \quad x_k = d_k \zeta, 
	\\
	 \gamma_k &= \dfrac{d_k}{d}, \quad 
	\delta_k = \dfrac{b_k}{d}, \quad \varepsilon_k = \dfrac{a_k}{d}, \quad d = \max \{ d_1, d_2, \dots, d_N \},
	\\
	\kappa_0 &= \kappa_2 d, \quad \chi_l(d_l \tau) = d_l \varphi_l(\tau), \quad \chi_l'(d_l \tau) = \varphi_l'(\tau), \quad \alpha_{kl} = \alpha_k - \alpha_l.
\end{align*}
%
for \( k, l = 1, 2, \dots, N \).

The system \eqref{PopovEq461086} must be supplemented by the following constraints:
\begin{equation} \label{PopovEq370089} 
	\int_{-1}^{1} \varphi_k'(\tau)\, d\tau = 0 \quad (k = 1, 2, \dots, N),
\end{equation}
which follow from the requirement that the crack faces remain closed at their edges.

The system of integro-differential equations \eqref{PopovEq461086} can be solved numerically, as described in \cite{Popov15,Popov17}. However, to avoid solving a high-dimensional system when the number of cracks is large, an iterative method is proposed. This method consists of solving \( N \) independent integro-differential equations at each iteration step \( i \) (with \( i \ge 1 \)):

\begin{equation} \label{PopovEq143660} 
	\begin{aligned}
		\begin{multlined}
			\frac{1}{2\pi} \int_{-1}^{1} \varphi'_{(i)k}(\tau) \left[ \dfrac{1}{\tau - \zeta} + F_{kk}(\tau - \zeta) \right] d\tau  
			\\
			+ \frac{1}{2\pi} \int_{-1}^{1} \varphi_{(i)k}(\tau) \left[ -\gamma_k^2 \kappa_0^2 \ln |\tau - \zeta| + U_{kk}(\tau - \zeta) \right] d\tau 
			\\
			= f_k(\zeta) - \sum_{\substack{l=1 \\ l \ne k}}^N \Bigg[ 
			\frac{1}{2\pi} \int_{-1}^{1} \varphi'_{(i-1)l}(\tau) F_{kl}(\tau, \zeta)  d\tau 
			\\ 
			+ \frac{1}{2\pi} \int_{-1}^{1} \varphi_{(i-1)l}(\tau) U_{kl}(\tau, \zeta) d\tau 
			\Bigg];
		\end{multlined}
		\\ 
		\displaystyle\int_{-1}^{1} \varphi'_{(i)k}(\tau) d\tau = 0 \quad (k = 1, 2, \dots, N; \; i = 1, 2, \dots)
	\end{aligned} 
\end{equation} 

The zero approximation is defined as the solution of equations \eqref{PopovEq143660} with the right-hand side reduced to \( f_k(\zeta) \), that is, without the sum terms. This corresponds to the displacements caused by incident wave loading acting on each crack individually. The numerical solution of system \eqref{PopovEq143660} follows the same approach as described in \cite{Popov15,Popov17}.

To account for the singularity, the derivatives of the unknown functions are expressed as
\begin{equation} \label{PopovEq992786} 
	\varphi'_{(i)k}(\tau) = \frac{\psi_{(i)k}(\tau)}{\sqrt{1 - \tau^2}}.
\end{equation} 
The functions \( \psi_{(i)k}(\tau) \) are then approximated by interpolation polynomials:
\begin{equation} \label{PopovEq972705} 
	\psi_{(i)k}(\tau) = \sum_{m=1}^{n} \psi_{(i)km} \frac{T_n(\tau)}{(\tau - \tau_m) T_n'(\tau)}, \quad 
	\psi_{(i)km} = \psi_{(i)k}(\tau_m),
\end{equation}
where \( T_n(\tau) \) is the Chebyshev polynomial of the second kind, and \( \tau_m \) (\( m = 1, 2, \dots, n \)) are its roots.

To determine \( \psi_{(i)k}(\tau) \), a system of linear algebraic equations is derived from \eqref{PopovEq143660} using the method of mechanical quadratures, as described in \cite{Popov15, Popov14}. After solving this system, the approximate values of the SIFs at the \( i \)th iteration step are computed using the following formulas:
\begin{equation} \label{PopovEq553928} 
	\begin{aligned} 
		K_{3l}^{\pm} &= G \sqrt{2d_l} \, k_{3l}^{\pm}, \quad
		k_{3l}^{+} = \dfrac{1}{n \sqrt{2\gamma_l}} \displaystyle\sum_{m=1}^{n} (-1)^m \psi_{(i)lm} \cot\left(\dfrac{\beta_m}{2}\right),
		\\
		k_{3l}^{-} &= \dfrac{1}{n \sqrt{2\gamma_l}} \displaystyle\sum_{m=1}^{n} (-1)^m \psi_{(i)lm} \tan\left(\dfrac{\beta_m}{2}\right), \quad 		
	\end{aligned} 
\end{equation}
where \(\beta_m = (2m - 1)\pi/(2n)\); \( l = 1, 2, \dots, N \).

The proposed iterative method was numerically implemented to evaluate its practical convergence, the stability of the iterative process, and its capacity to capture the influence of wave interactions between cracks on the SIF. To compare results obtained using different approaches, a configuration consisting of three cracks of equal length \( 2d \) was considered. The cracks are positioned at the vertices of an equilateral triangle with side length \( 4d \), as illustrated in Fig.~\ref{PopovFig02}.

\begin{figure}[t]
	\sidecaption[t]
	\includegraphics[width=.5\textwidth]{31_Popov/PopovFig02.jpeg}
	\caption{System of three cracks.}
	\label{PopovFig02}
\end{figure}

\begin{figure}[t]
	\sidecaption[t]
	\includegraphics[width=.5\textwidth]{31_Popov/PopovFig03.jpeg}
	\caption{Absolute values of the SIFs.}
	\label{PopovFig03}
\end{figure}

It is assumed that plane longitudinal shear waves propagate along the positive direction of the \( Oy \) axis and interact with the cracks. In Fig.~\ref{PopovFig03}, we present plots showing the dependence of the modulus of the dimensionless SIF, \( |k_1| = |k_{31}^{+}| = |k_{31}^{-}| \), on the dimensionless wave number \( \kappa_0 = \kappa_2 d \) for the first crack.

Curve 1 represents the SIF values obtained by directly solving the system of integro-differential equations \eqref{PopovEq461086}. Curves 2 and 4 show the SIF values obtained using formulas \eqref{PopovEq553928} after performing the corresponding number of iterations. It can be observed that, after four iterations, the results obtained by both methods are nearly identical. Further iterations lead to complete agreement with the direct solution. In the low-frequency range (\( \kappa_0 \leq 3 \)), convergence is already achieved by the second iteration. In the high-frequency range, where resonant peaks occur, full agreement with curve 1 is observed after 6--8 iterations.

\begin{figure}[b]
	\sidecaption[b]
	\includegraphics[width=.5\textwidth]{31_Popov/PopovFig04.jpeg}
	\caption{System of seven cracks.}
	\label{PopovFig04}
\end{figure}

\begin{figure}[t]
	\sidecaption[t]
	\includegraphics[width=.5\textwidth]{31_Popov/PopovFig05.jpeg}
	\caption{Absolute values of the SIFs at the adjacent edges of the cracks.}
	\label{PopovFig05}
\end{figure}

Next, the interaction of the same type of waves with a system of seven cracks (Fig.~\ref{PopovFig04}) is analyzed. The calculated values of the dimensionless SIFs—\( |k_{31}^{+}| \), \( |k_{32}^{-}| \), \( |k_{33}^{-}| \), and \( |k_{34}^{-}| \)—at the adjacent edges of the defects are presented in the plots in Fig.~\ref{PopovFig05}. As the frequency increases, the frequency dependence of the SIFs becomes more complex due to the emergence of resonant peaks. Notably, the lowest SIF values are observed for the fourth crack, indicating a shielding (screening) effect by the surrounding cracks. Initially, a monotonic increase in the SIF is observed for all cracks, which later transitions into an oscillatory pattern as the frequency continues to rise.


\section{Determination of Stress Intensity Factors in a Cylindrical Body with Cracks under Harmonic Oscillation of Longitudinal Shear}

We consider an infinite elastic cylindrical body whose axis is aligned with the \( Oz \) axis. The contour of its cross-section in a plane perpendicular to the axis is described by an arbitrary closed smooth curve, given in polar coordinates by the equation
\begin{equation} \label{31.17)} 
	r = r_0 \psi(\varphi), \quad 0 \le \varphi < 2\pi,
\end{equation} 
where \( r_0 \) is the characteristic geometric size of the cross-section. The cylinder contains \( N \) tunnel cracks, each occupying a segment of length \( 2d_k \) in the cross-section, with centers located at points \( O_k(a_k, b_k) \) for \( k = 1, 2, \dots, N \). All cracks are entirely enclosed within the domain defined by the cross-sectional boundary (Fig.~\ref{PopovFig06}).

\begin{figure}[b]
	\sidecaption[b]
	\includegraphics[width=.45\textwidth]{31_Popov/PopovFig06.jpeg}
	\caption{Cross-section of the cylindrical body.}
	\label{PopovFig06}
\end{figure}

Longitudinal shear oscillations in the cylinder arise due to a harmonic self-equilibrated load applied on the outer boundary, given by \( GP(\varphi)e^{-i\omega t} \), and directed along the \( Oz \) axis. Here, \( G \) is the shear modulus, \( \omega \) is the oscillation frequency, and \( P(\varphi) \) is the prescribed dimensionless amplitude of the load. The time-dependent multiplier \( e^{-i\omega t} \), which represents harmonic variation, is omitted in subsequent expressions, and only amplitude values are considered.

Under these loading conditions, only the \( z \)-component of the displacement vector, \( w(r, \varphi) \), is nonzero. This displacement satisfies the Helmholtz equation \eqref{PopovEq625098}, where \( \Delta \) denotes the Laplace operator in the polar coordinate system centered at point \( O \) (Fig.~\ref{PopovFig06}). The traction boundary condition on the cylinder surface takes the form
\begin{equation} \label{PopovEq725160} 
	\tau_{\bar{n}z} \big(r_0 \psi(\varphi), \varphi \big) = GP(\varphi), \quad 0 \leq \varphi < 2\pi,
\end{equation}
where \( \bar{n} \) is the unit normal vector to the boundary.

To formulate boundary conditions on the cracks, a local Cartesian coordinate system \( O_k x_k y_k \) is introduced for each crack. The relationship between the global and local Cartesian systems, as well as transformations among local systems of different cracks, is given by formulas \eqref{PopovEq491412}. The transformation from polar coordinates \( (r, \varphi) \) to the local system \( O_k x_k y_k \) is defined as
\begin{equation} \label{PopovEq876755} 
	\begin{array}{rcl} 
		x_l &=& \big(r\cos \varphi - a_l\big) \cos \alpha_l + \big(r\sin \varphi - b_l\big) \sin \alpha_l,
		\\ 
		y_l &=& -\big(r\cos \varphi - a_l\big) \sin \alpha_l + \big(r\sin \varphi - b_l\big) \cos \alpha_l.
	\end{array}
\end{equation}

Let \( W_k(x_k, y_k) \) denote the displacement obtained after transforming from polar coordinates to the local Cartesian coordinates using \eqref{PopovEq876755}. Assuming that the crack surfaces are free of stress, the boundary condition
\begin{equation} \label{PopovEq616707} 
	\tau_{zy_k}(x_k, 0) = 0, \quad -d_k < x_k < d_k, \quad (k = 1, 2, \dots, N)
\end{equation}
must be satisfied along each crack.

The displacements \( W_k(x_k, y_k) \) on the crack surfaces also exhibit discontinuities characterized by the unknown jumps defined in \eqref{PopovEq759458}. Under the formulated conditions, the problem reduces to determining the stress distribution in the body and the SIFs for each crack.

To solve this boundary value problem, discontinuous solutions of equation \eqref{PopovEq625098} are constructed for each crack in its associated local coordinate system. These solutions, incorporating the jump conditions \eqref{PopovEq759458}, are given by formulas \eqref{PopovEq185484} and \eqref{PopovEq953094}.

In the polar coordinate system, the total displacement \( W(r, \varphi) \) is then represented as
\begin{equation} \label{PopovEq232370} 
	W(r, \varphi) = W_0(r, \varphi) + \sum_{l=1}^{N} W_l(r, \varphi).
\end{equation}
Here, the functions \( W_l(r, \varphi) \) are the transformed discontinuous solutions, obtained from the local coordinates and expressed in terms of polar coordinates as
\begin{multline*}
	W_l(r, \varphi) = W_{dl} \Big[ (r\cos \varphi - c_l)\cos \alpha_l + (r\sin \varphi - d_l)\sin \alpha_l,\,
	\\
	(r\sin \varphi - d_l)\cos \alpha_l - (r\cos \varphi - c_l)\sin \alpha_l \Big],
\end{multline*}
and \( W_0(r, \varphi) \) is an unknown regular solution of the Helmholtz equation \eqref{PopovEq625098}, introduced to ensure that the boundary condition \eqref{PopovEq725160} is satisfied by the total field \eqref{PopovEq232370}.

The unknown function \( W_0(r, \varphi) \) is sought as a linear combination of \( M \) particular solutions of equation \eqref{PopovEq625098}:
\begin{equation}\label{PopovEq659953}
	W_0(r, \varphi) = r_0 \sum_{i=1}^{M} A_i\, g_i(r, \varphi),
\end{equation}
where the basis functions \( g_i(r, \varphi) \) are defined as
\[
g_{2m-1}(r, \varphi) = J_{m-1}(\kappa_2 r)\cos[(m - 1)\varphi], \quad
g_{2m}(r, \varphi) = J_m(\kappa_2 r)\sin(m \varphi),
\]
with \( J_m(\cdot) \) denoting the Bessel function of the first kind. The functions $g_{i} \left(r,\varphi \right)$are linearly independent and form a completely closed system of functions in the domain of cross section \cite{Popov24} of the body. In formula \eqref{PopovEq659953} $A_{i} $ are unknown coefficients, the number $M$ is subsequently chosen for reasons of ensuring the given accuracy of the solution. The constructed representation \eqref{PopovEq659953} of the displacements is a solution of equation \eqref{PopovEq625098} for arbitrary jumps \eqref{PopovEq759458} and coefficients $A_{i} $. Their definition is based on satisfying boundary conditions at the cracks and the cylinder surface. To implement the boundary conditions on cracks \eqref{PopovEq616707} in the coordinate system associated with the $k$th crack, the displacements \eqref{PopovEq232370} are given in the form:
\begin{equation} \label{PopovEq317159} 
	W_{k} \left(x_{k} ,y_{k} \right)=W_{k0} \left(x_{k} ,y_{k} \right)+\sum _{l=1}^{N}W_{kl} \left(x_{k} ,y_{k} \right) , 
\end{equation} 
where displacements $W_{k0} \left(x_{k} ,y_{k} \right)$ are obtained from $W_{0} \left(r,\varphi \right)$ after coordinate transformations according to \eqref{PopovEq876755}, and $W_{kl} \left(x_{k} ,y_{k} \right)$ ($l=1,\dots,N$) as a result of substitution $x_{l} ,y_{l} $ in $W_{dl} \left(x_{l} ,y_{l} \right)$, according to the second formulas \eqref{PopovEq491412}. After substitution of expressions \eqref{PopovEq317159} for displacements $W_{k} \left(x_{k} ,y_{k} \right)$ in the boundary condition \eqref{PopovEq616707} system of integro-differential equations same as \eqref{PopovEq461086} is obtained. In this system
\begin{equation} \label{PopovEq793073} 
	\begin{array}{l} 
		\chi _{l} \left(d_{l} \tau \right)=d_{l} \varphi _{l} \left(\tau \right);\; \chi '_{l} \left(d_{l} \tau \right)=\varphi '_{l} \left(\tau \right); 
		\\
		f_{k} (\zeta )=\displaystyle\sum _{i=1}^{M}A_{i} f{}_{ik}  (\zeta ),\; f_{ik} (\zeta )=-\dfrac{\partial g_{i} (r_{0k} ,\varphi _{0k} )}{\partial y_{k} } 
		\\
		r_{0k}^{2} =(a_{0k} +\gamma _{k} \zeta \cos \alpha _{k} )^{2} +(b_{0k} +\gamma _{k} \zeta \sin \alpha _{k} )^{2} ,
		\\
		\tan\varphi _{0k} =\dfrac{b_{0k} +\gamma _{k} \zeta \sin \alpha _{k} }{a_{0k} +\gamma _{k} \zeta \cos \alpha _{k} } ;\; a_{0k} =\dfrac{a_{k} }{r_{0} } ,\; b_{0k} =\dfrac{b_{k} }{r_{0} } ,\;\gamma _{k} =\dfrac{d_{k} }{r_{0} } .
	\end{array} 
\end{equation} 
Now, if the unknown functions \( \varphi_l(\tau) \) from \eqref{PopovEq793073} are represented as a linear combination of new unknown functions,
\begin{equation} \label{PopovEq416497} 
	\varphi_l(\tau) = \sum_{i=1}^{M} A_i \varphi_{il}(\tau), \quad 
	\varphi_l'(\tau) = \sum_{i=1}^{M} A_i \left(\varphi_{il}(\tau)\right)', 
\end{equation}
then the auxiliary functions \( \varphi_{il}(\tau) \) \((i=1,2,\dots,M;\; l=1,2,\dots,N)\) can be obtained by solving \( M \) systems of equations of the form \eqref{PopovEq461086} with the corresponding right-hand sides \( f_{ik}(\zeta) \). 

The solution of the integral equations \eqref{PopovEq461086} and \eqref{PopovEq370089} relies on the approximation of the derivatives of the unknown functions using formulas \eqref{PopovEq992786} and \eqref{PopovEq972705}, together with the application of the mechanical quadrature method. This approach reduces the problem to solving systems of linear algebraic equations for the coefficients \( \psi_{ikm} = \psi_{ik}(\tau_m) \), which differ only in their right-hand sides, while the system matrices remain the same.

Once the functions \( \varphi_{il}(\tau) \) are determined, the boundary conditions on the crack surfaces are automatically satisfied for any values of the coefficients \( A_i \). To complete the solution of the formulated problem, it is necessary to determine these coefficients from the boundary condition \eqref{PopovEq725160} on the external boundary of the body. 

For this purpose, the boundary shear stress is computed as
\begin{equation} \label{PopovEq749104} 
	\tau_{\bar{n}z}(r_0 \psi(\varphi), \varphi) = \tau_{xz}(r_0 \psi(\varphi), \varphi) A_x + \tau_{yz}(r_0 \psi(\varphi), \varphi) A_y,
\end{equation}
where \( \bar{n} \) is the unit normal vector to the surface, and \( A_x \), \( A_y \) are the direction cosines of the vector \( \bar{n} \) with respect to the Cartesian axes.

 After substituting the expressions for the stresses from \eqref{PopovEq749104} into the boundary condition \eqref{PopovEq725160}, and applying the representations \eqref{PopovEq416497}, the boundary condition takes the form:
\begin{equation} \label{PopovEq671340} 
	\sum_{i=1}^{M} A_i \left( \sum_{l=1}^{N} \int_{-1}^{1} \varphi_{il}(\tau) G_l(\tau, \varphi) \, d\tau + B_i(\varphi) \right) = P(\varphi),
\end{equation}
where

\begin{align*}
	G_l(\tau, \varphi) &= \dfrac{i \kappa_0^2 \gamma_l^2}{4} \left( E_{1l}(\tau, \varphi)\, A_x + E_{2l}(\tau, \varphi)\, A_y \right),
	\\
	B_i(\varphi) &= f_{1i}(\varphi)\, A_x + f_{2i}(\varphi)\, A_y,
\end{align*}
%
with the auxiliary functions defined as:

\begin{align*}
	E_1(\tau, \varphi)  &=
	\begin{multlined}[t]
		\dfrac{-y_0\left((\gamma \tau - x_0)\cos \alpha + y_0 \sin \alpha \right)}{R_0^2} H_0^{(1)}(\kappa_0 R_0)
		\\
		+ \dfrac{H_1^{(1)}(\kappa_0 R_0)}{\kappa_0 R_0^3} \cdot \sin \alpha \left( y_0^2 - (\gamma \tau - x_0)^2 \right) + 2 y_0 (\gamma \tau - x_0) \cos \alpha,
	\end{multlined}
	\\
	E_2(\tau, \varphi) &= 
	\begin{multlined}[t]
		\dfrac{-y_0\left((\gamma \tau - x_0)\sin \alpha - y_0 \cos \alpha \right)}{R_0^2} H_0^{(1)}(\kappa_0 R_0)
		\\
		+ \dfrac{H_1^{(1)}(\kappa_0 R_0)}{\kappa_0 R_0^3} \cdot \cos \alpha \left( -y_0^2 + (\gamma \tau - x_0)^2 \right) + 2 y_0 (\gamma \tau - x_0) \sin \alpha,
	\end{multlined}
\end{align*}
%
and
\[
\begin{array}{rl}
	R_0 =& \sqrt{(\gamma \tau - x_0)^2 + y_0^2}, \quad
	f_{1i} = r_0 \dfrac{\partial g_i}{\partial x}, \quad
	f_{2i} = r_0 \dfrac{\partial g_i}{\partial y}, \\[0.5em]
	x_0 =& \psi(\varphi)\cos(\varphi - \alpha) - a_0 \cos \alpha - b_0 \sin \alpha, \\[0.5em]
	y_0 =& \psi(\varphi)\sin(\varphi - \alpha) - b_0 \cos \alpha + a_0 \sin \alpha.
\end{array}
\]
The coefficients \( A_i \) are determined approximately from the system of linear algebraic equations obtained by applying the collocation method to equation \eqref{PopovEq671340}. The collocation points are chosen as \( \sigma_s = 2\pi s / M \) for \( s = 1, \dots, M \). 

Once the integral equations \eqref{PopovEq143660} have been solved and the coefficients \( A_i \) determined, the approximate values of the SIFs for each crack are computed using the following formulas:
\begin{equation} \label{31.28)} 
	K_{3l}^{\pm} = G \sqrt{2d_l} k_{3l}^{\pm},
\end{equation}
where

\begin{align*} 
	k_{3l}^{+} &= \dfrac{1}{2n} \sum_{i=1}^{M} A_i \displaystyle\sum_{m=1}^{n} (-1)^m \psi_{ilm} \cot \dfrac{\beta_m}{2},
	\\
	k_{3l}^{-} &= \dfrac{(-1)^{n+1}}{2n} \displaystyle\sum_{i=1}^{M} A_i \sum_{m=1}^{n} (-1)^m \psi_{ilm} \tan \dfrac{\beta_m}{2}.
\end{align*}



As an example of the numerical implementation of the proposed method, an elliptical cylinder containing four identical strip-like cracks arranged as shown in Fig.~\ref{PopovFig07} was considered. The ellipse was assumed to have eccentricity \( \varepsilon = 0.5 \), and the inclination of the cracks relative to the major axis was characterized by the angle \( \alpha \) (see Fig.~\ref{PopovFig07}). The load applied to the lateral surface of the cylinder was specified by the expression \( P(\varphi) = \sin 2\varphi \). Due to the symmetry of the loading and the cracks' arrangement, the behavior of the SIF can be studied for a single representative crack, for instance, the one located in the first quadrant of the cross section.

\begin{figure}[t]
	\sidecaption[t]
	\includegraphics[width=.45\textwidth]{31_Popov/PopovFig07.png}
	\caption{Cross section of the elliptical cylinder with four cracks.}
	\label{PopovFig07}
\end{figure}



The principal varying parameter was \( \gamma = d / r_0 \), representing the ratio of crack half-length to the semi-major axis of the ellipse. For the given ellipse and at an angle \( \alpha = 30^\circ \), the crack intersects the boundary of the cylinder when \( \gamma \approx 0.55 \). The calculation results are presented in the form of plots showing the dependence of the dimensionless SIF on the dimensionless frequency \( \kappa_0 = \kappa_2 r_0 \). Figure~\ref{PopovFig08} illustrates the influence of the relative crack length on the SIF for values of \( \gamma \) in the range \( 0 < \gamma < 0.52 \). Curves 1, 2, and 3 correspond to \( \gamma = 0.2 \), \( 0.4 \), and \( 0.52 \), respectively.

\begin{figure}[t]
	\sidecaption[t]
	\includegraphics[width=.5\textwidth]{31_Popov/PopovFig08.jpg}
	\caption{Absolute values of the SIF around the edge of crack.}
	\label{PopovFig08}
\end{figure}

The results indicate that the crack length significantly affects the resonance frequency: as the crack length increases and the tips approach the boundary (i.e., increasing \( \gamma \)), the resonance frequency decreases. When the crack length approaches the threshold \( 0.52 < \gamma < 0.55 \), such that the cracks are nearly intersecting the boundary, the effectiveness of the proposed method deteriorates.


\section{Determination of the Stress State and Stress Intensity Factors Under Dynamic Loading of a System of Cracks}

Consider an isotropic, unbounded elastic medium under plane strain conditions, containing through-cracks arranged as shown in Fig.~\ref{PopovFig01}. Normal ($N_{k} e^{-i\omega t}$) or shear ($T_{k} e^{-i\omega t}$) self-equilibrated loadings are applied to the crack surfaces.

Let \( u(x,y) \) and \( v(x,y) \) denote the displacement components of the wave field generated by the applied loading. These displacements satisfy the equations of motion for a plane strain state:
\begin{equation} \label{PopovEq842557} 
	\begin{aligned}
		\left(\lambda +2\mu \right)\dfrac{\partial }{\partial x} \left(\dfrac{\partial u}{\partial x} +\dfrac{\partial v}{\partial y} \right)+\mu \Delta u&=-\rho \omega ^{2} u,
		\\[0.5em]
		\left(\lambda +2\mu \right)\dfrac{\partial }{\partial y} \left(\dfrac{\partial u}{\partial x} +\dfrac{\partial v}{\partial y} \right)+\mu \Delta v&=-\rho \omega ^{2} v, 
	\end{aligned}
\end{equation} 
where \( \lambda \) and \( \mu \) are the Lamé constants, and \( \rho \) is the density of the elastic medium.

To formulate the boundary conditions on the crack surfaces, we associate with each crack a local coordinate system \( O_{k}x_{k}y_{k} \) (see Fig.~\ref{PopovFig01}). Let \( u_{k}(x_{k}, y_{k}) \), \( v_{k}(x_{k}, y_{k}) \), \( \tau_{kxy}(x_{k}, y_{k}) \), \( \sigma_{kx}(x_{k}, y_{k}) \), and \( \sigma_{ky}(x_{k}, y_{k}) \) denote the displacements and stress components in the coordinate system associated with the $k$th crack. Then, on the loaded surfaces of the cracks, the following boundary conditions are imposed:
\begin{equation} \label{PopovEq491114} 
	\sigma_{ky}(x_{k}, 0) = N_{k},\quad \tau_{kxy}(x_{k}, 0) = T_{k},\quad -d_{k} < x_{k} < d_{k}
\end{equation} 
for \( k = 1, 2, \dots, N \).

In addition, displacement discontinuities occur on the crack surfaces. The unknown jumps in displacements are denoted by:
\begin{equation} \label{PopovEq200550} 
	\begin{array}{rcl}
		v_{k}(x_{k}, +0) - v_{k}(x_{k}, -0) &=& \chi_{3k}(x_{k}),
		\\[0.5em]
		u_{k}(x_{k}, +0) - u_{k}(x_{k}, -0) &=& \chi_{4k}(x_{k}),\quad -d_{k} < x_{k} < d_{k}.  
	\end{array}
\end{equation} 

Under such conditions, it is necessary to determine the displacements and stresses in the body with defects and to derive formulas for calculating the SIFs. Following the approach in \cite{Popov14}, the solution is based on the use of discontinuous solutions to the equations \eqref{PopovEq842557}, constructed for each defect in its associated local coordinate system \( O_{l}x_{l}y_{l} \) ($l = 1, 2, \dots, N$). For cracks, these discontinuous solutions exhibit displacement jumps \eqref{PopovEq759458} and are given by the expressions \cite{Popov16}:

\begin{equation} \label{PopovEq144429} 
	\begin{aligned}
		v_{dl}(x_{l}, y_{l}) &=
		\begin{multlined}[t]
			\int_{-d_{l}}^{d_{l}} \chi_{3l}(\eta)\, G_{33}(\eta - x_{l}, y_{l})\, d\eta 
			\\[-1em]
			+ \int_{-d_{l}}^{d_{l}} \chi_{4l}(\eta)\, G_{34}(\eta - x_{l}, y_{l})\, d\eta,
		\end{multlined}
		\\
		u_{dl}(x_{l}, y_{l}) &=
		\begin{multlined}[t]
			\int_{-d_{l}}^{d_{l}} \chi_{3l}(\eta)\, G_{43}(\eta - x_{l}, y_{l})\, d\eta 
			\\
			+ \int_{-d_{l}}^{d_{l}} \chi_{4l}(\eta)\, G_{44}(\eta - x_{l}, y_{l})\, d\eta.
		\end{multlined}
	\end{aligned}
\end{equation}

Here,
\[
\begin{array}{rcl}
	G_{33} &=& \frac{1}{\kappa_2^2} \frac{\partial}{\partial y_{l}} \left[\left(\kappa_2^2 + 2\frac{\partial^2}{\partial x_{l}^2}\right) r_1 - 2 \frac{\partial^2 r_2}{\partial x_{l}^2} \right],
	\\[0.5em]
	G_{34} &=& \frac{1}{\kappa_2^2} \frac{\partial}{\partial y_{l}} \left[ 2\left(\kappa_1^2 + 2\frac{\partial^2}{\partial x_{l}^2} \right) r_1 - \left(\kappa_2^2 + 2\frac{\partial^2}{\partial x_{l}^2} \right) r_2 \right],
	\\[0.5em]
	G_{43} &=& \frac{1}{\kappa_2^2} \frac{\partial}{\partial y_{l}} \left[ \left(\kappa_2^2 + 2\frac{\partial^2}{\partial x_{l}^2} \right) r_1 - 2\left(\kappa_2^2 + \frac{\partial^2}{\partial x_{l}^2} \right) r_2 \right],
	\\[0.5em]
	G_{44} &=& \frac{1}{\kappa_2^2} \frac{\partial}{\partial y_{l}} \left[ -2\frac{\partial^2 r_1}{\partial x_{l}^2} + \left(\kappa_2^2 + 2\frac{\partial^2}{\partial x_{l}^2} \right) r_2 \right],
\end{array}
\]
with
\begin{align*}
&r_j(\eta - x_{l}, y_{l}) = -\frac{i}{4} H_0^{(1)}\left[\kappa_j \sqrt{(\eta - x_{l})^2 + y_{l}^2} \right] \quad (j = 1,2),
\\
&\kappa_1^2 = \frac{\omega^2 \rho}{\lambda + 2\mu}, \quad \kappa_2^2 = \frac{\omega^2 \rho}{\mu}.
\end{align*}

The corresponding stress components are expressed as:

\begin{equation} \label{PopovEq460990} 
	\begin{aligned} 
		\sigma_{dx_{l}}(x_{l}, y_{l}) &=
		\begin{multlined}[t]
			\mu \int_{-d_{l}}^{d_{l}} \chi_{3l}'(\eta)\, E_{03}(\eta - x_{l}, y_{l})\, d\eta 
			\\
			+ \mu (\kappa_2^2 - 2\kappa_1^2) \int_{-d_{l}}^{d_{l}} \chi_{3l}(\eta)\, r_1(\eta - x_{l}, y_{l})\, d\eta 
			\\
			+ \mu \int_{-d_{l}}^{d_{l}} \chi_{4l}'(\eta)\, E_{04}(\eta - x_{l}, y_{l})\, d\eta,
		\end{multlined}
		\\[0.5em]
		\sigma_{dy_{l}}(x_{l}, y_{l}) &=
		\begin{multlined}[t]
			\mu \int_{-d_{l}}^{d_{l}} \chi_{3l}'(\eta)\, E_{13}(\eta - x_{l}, y_{l})\, d\eta 
			\\
			+ \mu \int_{-d_{l}}^{d_{l}} \chi_{4l}'(\eta)\, E_{14}(\eta - x_{l}, y_{l})\, d\eta 
			\\
			- \mu \kappa_2^2 \int_{-d_{l}}^{d_{l}} \chi_{3l}(\eta)\, r_1(\eta - x_{l}, y_{l})\, d\eta,
		\end{multlined}
		\\[0.5em]
		\tau_{dy_{l} x_{l}}(x_{l}, y_{l}) &=
		\begin{multlined}[t]
			\mu \int_{-d_{l}}^{d_{l}} \chi_{3l}'(\eta)\, E_{23}(\eta - x_{l}, y_{l})\, d\eta 
			\\
			+ \mu \int_{-d_{l}}^{d_{l}} \chi_{4l}'(\eta)\, E_{24}(\eta - x_{l}, y_{l})\, d\eta 
			\\
			+ \mu \kappa_2^2 \int_{-d_{l}}^{d_{l}} \chi_{3l}(\eta)\, r_2(\eta - x_{l}, y_{l})\, d\eta.
		\end{multlined}
	\end{aligned}
\end{equation}

When deriving the expressions in \eqref{PopovEq460990}, integration by parts was applied to reduce the singularity order of the integrable functions. This was done under the assumption that the displacement jumps vanish at the crack tips, i.e., \( \chi_{3l}(\pm d_{l}) = \chi_{4l}(\pm d_{l}) = 0 \). The following notations are used for compactness:


\begin{align*}
	E_{03} &= E_{24} = -\dfrac{4}{\kappa_2^2} \left[ \left( \kappa_1^2 + \dfrac{\partial^2}{\partial \eta^2} \right) \dfrac{\partial r_1}{\partial \eta} - \left( \kappa_2^2 + \dfrac{\partial^2}{\partial \eta^2} \right) \dfrac{\partial r_2}{\partial \eta} \right],
	\\
	E_{04} &= -\dfrac{2}{\kappa_2^2} \left[ \left( 2\dfrac{\partial^2}{\partial \eta^2} + 2\kappa_1^2 - \kappa_2^2 \right) \dfrac{\partial r_1}{\partial y_l} - \left( 2\dfrac{\partial^2}{\partial \eta^2} + \kappa_2^2 \right) \dfrac{\partial r_2}{\partial y_l} \right],
	\\
	E_{13} &= \dfrac{4}{\kappa_2^2} \left[ \left( \kappa_2^2 + \dfrac{\partial^2}{\partial \eta^2} \right) \dfrac{\partial r_1}{\partial \eta} - \left( \kappa_2^2 + \dfrac{\partial^2}{\partial \eta^2} \right) \dfrac{\partial r_2}{\partial \eta} \right],
	\\
	E_{14} &= E_{23} = \dfrac{2}{\kappa_2^2} \left[ \left( \kappa_2^2 + 2\dfrac{\partial^2}{\partial \eta^2} \right) \dfrac{\partial r_1}{\partial y_l} - \left( \kappa_2^2 + 2\dfrac{\partial^2}{\partial \eta^2} \right) \dfrac{\partial r_2}{\partial y_l} \right].
\end{align*}

The following expressions transform the displacements and stresses of the discontinuous solutions \eqref{PopovEq144429} and \eqref{PopovEq460990} from the local coordinate system \( O_l x_l y_l \) to the global system \( Oxy \):

\begin{equation} \label{PopovEq642436} 
	\begin{array}{rcl}
		u_{gl} &=& u_{dl} \cos \alpha_l - v_{dl} \sin \alpha_l,
		\\
		v_{gl} &=& u_{dl} \sin \alpha_l + v_{dl} \cos \alpha_l,
		\\
		\sigma_{glx} &=& \sigma_{dx_l} \cos^2 \alpha_l + \sigma_{dy_l} \sin^2 \alpha_l - \tau_{dy_l x_l} \sin 2\alpha_l,
		\\
		\sigma_{gly} &=& \sigma_{dx_l} \sin^2 \alpha_l + \sigma_{dy_l} \cos^2 \alpha_l + \tau_{dy_l x_l} \sin 2\alpha_l,
		\\
		2\tau_{glyx} &=& \sigma_{dx_l} \sin 2\alpha_l + \sigma_{dy_l} \sin 2\alpha_l + 2\tau_{dy_l x_l} \cos 2\alpha_l.
	\end{array}
\end{equation}

In formulas \eqref{PopovEq642436} and further, $\alpha _{l} $ is the angle between the axis $Ox$ and the axis $O_{l} x_{l} $. Then the displacements wave field can be represented as
\begin{equation} \label{PopovEq808654} 
	u(x,y)=\sum _{l=1}^{N}u_{gl} (x,y),\quad v(x,y)=\sum _{l=1}^{n}v_{gl} (x,y) . 
\end{equation} 
Stresses $\sigma _{x}$, $\sigma _{y}$, $\tau _{yx}$ are also determined by similar formulas. Thus, according to formulas \eqref{PopovEq808654}, displacements and stress in the body will be determined, provided that the unknown jumps \eqref{PopovEq200550} are determined. For this conditions \eqref{PopovEq491114} should be used. The stresses that enter these conditions, through discontinuous solutions \eqref{PopovEq144429}, \eqref{PopovEq460990} are expressed by the formulas
\begin{equation} \label{PopovEq347035} 
	\sigma _{ky} (x_{k} ,y_{k} )=\sum _{l=1}^{N}\sigma _{kl} (x_{k} ,y_{k} ),\quad \tau _{kyx} (x_{k} ,y_{k} )=\sum _{l=1}^{N}\tau _{kl} (x_{k} ,y_{k} ) ,     
\end{equation} 
\[
\begin{array}{rcl} \sigma _{kl}&=&\sigma _{dx_{l} } \sin ^{2} \alpha _{kl} +\sigma _{dy_{l} } \cos ^{2} \alpha _{kl} -\tau _{dy_{l} x_{l} } \sin 2\alpha _{kl}, 
	\\
	2\tau _{kl} &=&-\sigma _{dx_{l} } \sin 2\alpha _{kl} +\sigma _{dy_{l} } \sin 2\alpha _{kl} +2\tau _{dy_{l} x_{l} } \cos 2\alpha _{kl},\quad\alpha _{kl} =\alpha _{k} -\alpha _{l}. 
\end{array}
\] 
Substituting expressions \eqref{PopovEq347035} into the boundary conditions \eqref{PopovEq491114} with respect to the displacement jumps leads to a system of singular integro-differential equations. After isolating the singular terms and performing several transformations, the system takes the form:

\begin{equation} \label{PopovEq123113} 
	\begin{aligned}
		\begin{multlined}[b]
			\frac{1}{2\pi } \int _{-1}^{1}\varphi '_{3k} (\tau ) \left[-\frac{2\left(1-\xi ^{2} \right)}{\tau -\xi } +F_{13kk} (\tau -\xi )\right]d\tau 
			\\
			+\frac{1}{2\pi } \int _{-1}^{1}\varphi _{3k} (\tau ) \Big[\kappa _{0}^{2} \gamma _{k}^{2} \ln \left|\tau -\varsigma \right|+U_{13kk} (\tau -\varsigma)\Big]d\tau +
			\\
			+\frac{1}{2\pi }\sum _{\substack{l=1 \\l\ne k}}^{N}\Bigg[ \int _{-1}^{1}\varphi '_{3l} (\tau )F_{13kl} (\tau ,\varsigma )d\tau 
			+ \int _{-1}^{1}\varphi '_{4k} (\tau )F_{14kl} (\tau ,\varsigma ) d\tau  \Bigg] 
			\\
			+\frac{1}{2\pi }\sum_{\substack{l=1 \\ l\ne k}}^{N}\Bigg[ \int _{-1}^{1}\varphi _{3l} (\tau )U_{13kl} (\tau ,\varsigma )d\tau 
			\\
			+ \int _{-1}^{1}\varphi _{4k} (\tau )U_{14kl} (\tau ,\varsigma ) d\tau  \Bigg] 
		\end{multlined}&=N_{0k}  
		\\
		\begin{multlined}[b]
			\frac{1}{2\pi } \int_{-1}^{1}\varphi '_{4k} (\tau ) \left[\frac{2\left(1-\xi ^{2} \right)}{\tau -\varsigma } +F_{24kk} \left(\tau -\varsigma \right)\right]d\tau
			\\
			+\frac{1}{2\pi } \int_{-1}^{1}\varphi _{4k} (\tau ) \left[\kappa _{0}^{2} \gamma _{k}^{2} \ln \left|\tau -\varsigma \right|+U_{24kk} (\tau -\varsigma)\right]d\tau 
			\\
			+\frac{1}{2\pi }\sum_{\substack{l=1 \\l\ne k}}^{N}\Bigg[ \int _{-1}^{1}\varphi '_{3l} (\tau )F_{23kl} (\tau ,\varsigma )d\tau 
			\\
			+ \int _{-1}^{1}\varphi '_{4k} (\tau )F_{24kl} (\tau ,\varsigma ) d\tau  \Bigg] 
			\\
			+\frac{1}{2\pi }\sum_{\substack{l=1 \\l\ne k}}^{N}\Bigg[ \int _{-1}^{1}\varphi _{3l} (\tau )U_{23kl} (\tau ,\varsigma )d\tau 
			\\
			+\int _{-1}^{1}\varphi _{4k} (\tau )U_{24kl} (\tau ,\varsigma ) d\tau  \Bigg] 
		\end{multlined}&=T_{0k} 
	\end{aligned}
\end{equation} 
for \( k = 1, 2, \dots, N \). The functions \( F_{kl1s}, F_{kl2s}, U_{kl1s}, U_{kl2s} \) with \( s = 3, 4 \) are continuous in the domain \( -1 \le \varsigma, \tau \le 1 \). The following notations are used:

\[
\begin{array}{l}
	\varphi _{sk} (\tau )=d_{k}^{-1} \chi _{sk} (d_{k} \tau ),\;\; \gamma _{k} =d^{-1} d_{k} ,\;\; N_{0k} =\mu ^{-1} N_{k} ;\;\; T_{0k} =\mu ^{-1} T_{k}
	\\
	(k=1,2,\dots,N;\; s=3,4)
	\\
	d=\max (d_{1} ,d_{2} ,\dots,d_{N} );\;\; \kappa _{0} =\kappa _{2} d;\;\; \xi ^{2} =\dfrac{\kappa _{2}^{2} }{\kappa _{1}^{2} } =\dfrac{1-2\nu }{2(1-\nu )} 
\end{array}
\]
where \( \nu \) is the Poisson’s ratio of the elastic medium.

\noindent Additionally, the following conditions, expressing the smooth closure of crack faces, must be imposed:
\[
\int _{-1}^{1}\varphi '_{sk} (\tau )d\tau =0\quad (k=1, 2,\dots,N;\; s=3,4).
\] 

To solve system \eqref{PopovEq123113}, an iterative method is proposed, similar to that used for system \eqref{PopovEq461086}. At the \( i \)-th iteration step (\( i \ge 1 \)), the following system of \( 2N \) integro-differential equations is solved:

\begin{equation} \label{PopovEq994019} 
	\begin{aligned}
		\begin{multlined}[t]
			\frac{1}{2\pi} \int_{-1}^{1}\varphi'_{(i)3k}(\tau ) \Bigg[-\frac{2\left(1-\xi ^{2} \right)}{\tau -\xi} +F_{13kk} (\tau -\xi )\Bigg]d\tau 
			\\
			+\frac{1}{2\pi } \int_{-1}^{1}\varphi_{(i)3k} (\tau ) \left[\kappa _{0}^{2} \gamma _{k}^{2} \ln |\tau -\varsigma |+U_{13kk} \left(\tau -\varsigma \right)\right]d\tau 
			\\
			=N_{0k} -\frac{1}{2\pi }\sum _{\substack{l=1 \\ l\ne k}}^{N}\Bigg[ \int _{-1}^{1}\varphi '_{(i-1)3l} (\tau )F_{13kl} (\tau ,\varsigma )d\tau 
			\\
			+ \int _{-1}^{1}\varphi '_{(i-1)4k} (\tau )F_{14kl} (\tau ,\varsigma ) d\tau  \Bigg] 
			\\
			-\frac{1}{2\pi }\sum _{\substack{l=1 \\ l\ne k}}^{N}\Bigg[ \int _{-1}^{1}\varphi_{(i-1)3l} (\tau )U_{kl}^{23} (\tau ,\varsigma )d\tau 
			\\
			+ \int _{-1}^{1}\varphi _{(i-1)4l} (\tau )U_{kl}^{23} (\tau ,\varsigma ) d\tau  \Bigg],  
		\end{multlined}
		\\
		\begin{multlined}[b]
			\frac{1}{2\pi } \int _{-1}^{1}\varphi '_{(i)4k} (\tau ) \left[\frac{2\left(1-\xi ^{2} \right)}{\tau -\varsigma } +F_{24kk} \left(\tau -\varsigma \right)\right]d\tau 
			\\
			+\frac{1}{2\pi } \int _{-1}^{1}\varphi _{(i)4k} (\tau ) \left[\kappa _{0}^{2} \gamma _{k}^{2} \ln \left|\tau -\varsigma \right|+U_{24kk} \left(\tau -\varsigma \right)\right]d\tau =
			\\
			=T_{0k} -\frac{1}{2\pi }\sum_{\substack{l=1 \\ l\ne k}}^{N}\Bigg[ \int _{-1}^{1}\varphi '_{(i-1)3l} (\tau )F_{23kl} (\tau ,\varsigma )d\tau 
			\\
			+\int _{-1}^{1}\varphi '_{(i-1)4k} (\tau )F_{24kl} (\tau ,\varsigma ) d\tau  \Bigg] 
			\\
			-\frac{1}{2\pi }\sum_{\substack{l=1 \\ l\ne k}}^{N}\Bigg[ \int _{-1}^{1}\varphi _{(i-1)3l} (\tau )U_{23kl} (\tau ,\varsigma )d\tau 
			\\
			+\int _{-1}^{1}\varphi _{(i-1)4k} (\tau )U_{24kl} (\tau ,\varsigma ) d\tau  \Bigg]
		\end{multlined}
	\end{aligned}
\end{equation} 
for \( k = 1, 2, \dots, N \), \( i = 1, 2, \dots \).

The solutions of equations \eqref{PopovEq994019} with right-hand sides reduced to $N_{0k}$ and $T_{0k}$ (i.e., without interaction terms) are used as the zero approximation: \(\varphi_{(0)3k}(\tau)\), \(\varphi_{(0)4k}(\tau)\). In other words, the initial approximation corresponds to the solutions for \(N\) independent single cracks subjected to dynamic loading.

The numerical solution of system \eqref{PopovEq994019} is performed analogously to the solution of system \eqref{PopovEq143660}. The derivatives of the unknown functions, \(\varphi'_{(i)sk}(\tau)\), accounting for their singular behavior, are approximated using the expressions \eqref{PopovEq992786} and \eqref{PopovEq972705}. 

To compute the nodal values \(\psi_{(i)skm} = \psi_{(i)sk}(\tau_m)\), a system of linear algebraic equations is derived from \eqref{PopovEq994019} using the method of mechanical quadratures. After solving this system, the approximate dimensionless values of the tearing and shear SIF at the crack tips for the \(i\)th iteration are determined using the formulas from \cite{Popov14}:

\begin{equation} \label{PopovEq244046}
	\begin{array}{rcl}
		K^{\pm}_{(s-2)l} &=& \mu \sqrt{d_l}  k^{\pm}_{(s-2)l}, 
		\\[0.5em]
		k^{+}_{(s-2)l} &=& -\dfrac{1 - \xi^2}{n} \displaystyle\sum_{m=1}^{n} (-1)^m \psi_{(i)slm} \cot\left( \dfrac{\beta_m}{2} \right),
		\\[0.5em]
		k^{-}_{(s-2)l} &=& -\dfrac{1 - \xi^2}{n} (-1)^{s+n} \displaystyle\sum_{m=1}^{n} (-1)^m \psi_{(i)slm} \tan\left( \dfrac{\beta_m}{2} \right),
	\end{array}
\end{equation}
where \( l = 1, 2, \dots, N \), and \( s = 3, 4 \).

The numerical implementation of the proposed iterative method was conducted to assess its practical convergence, the stability of the iterative process, and its effectiveness in evaluating the influence of wave interactions between cracks on the SIFs. 

To compare the results obtained by directly solving system \eqref{PopovEq123113}—as was done in \cite{Popov14}—with those derived from the iterative method, a configuration consisting of four cracks of equal length \(2d\) was analyzed (see Fig.~\ref{PopovFig09}). The geometric arrangement of this crack system is characterized by the parameters \(\theta\), \(h\), and \(\beta\). All cracks are subjected to identical normal (\(N_{0k} = 1\)) or shear (\(T_{0k} = 1\)) loading conditions for each \(k=1,2,3,4\). Under these symmetric conditions, the SIFs at the tips of all cracks are expected to be identical.

\begin{figure}[b]
	\sidecaption[b]
	\includegraphics[width=.4\textwidth]{31_Popov/PopovFig09.jpg}
	\caption{Configuration of a system of four cracks of equal length.}
	\label{PopovFig09}
\end{figure}

The results of the iterative process are illustrated in Fig.~\ref{PopovFig10}, which shows the dependence of the computed dimensionless SIFs on the dimensionless frequency for different numbers of iterations. These results validate the convergence of the proposed method and its applicability for analyzing the collective dynamic behavior of interacting cracks.

\begin{figure}[t]
	\sidecaption[t]
	\includegraphics[width=.45\textwidth]{31_Popov/PopovFig10.jpg}
	\caption{Relationship between SIF and dimensionless frequency for different numbers of iterations.}
	\label{PopovFig10}
\end{figure}


The graphs in Fig.~\ref{PopovFig10} illustrate the variation in the absolute value of the SIF for normal stresses as a function of the dimensionless wave number \(\kappa_0 = \kappa_2 d\). The calculations were performed using formulas \eqref{PopovEq244046} for the parameters \(\theta = h = 0.75d\), \(\beta = 30^\circ\). The number of iterations ranged from 2 to 16. The dashed curve represents the SIF values for a single isolated crack, while the curve labeled “0” corresponds to the results obtained via direct solution of the integro-differential system \eqref{PopovEq123113}. The remaining curves show the SIF values calculated after the specified number of iterations.

As \(\kappa_0 \to 0\), the SIF approaches its static counterpart. The comparison reveals that the iterative method converges well: results from different methods align closely after the fourth iteration, with the exception of the first and most pronounced peak, which aligns only after 16 iterations. Notably, the extreme SIF values significantly exceed both their static analogues and the values for a single isolated crack.

However, the convergence of the iterative method degrades as \(\beta \to 0\), i.e., when the system contains closely spaced, nearly parallel cracks. In such cases, convergence is generally achieved only when \(h > 0.5d\). The results of convergence studies for \(\beta = 0\) and \(\theta = h = 0.75d\) are presented in Fig.~\ref{PopovFig11}.

\begin{figure}[b]
	\sidecaption[b]
	\includegraphics[width=.45\textwidth]{31_Popov/PopovFig11.jpg}
	\caption{Convergence for a system of four parallel cracks.}
	\label{PopovFig11}
\end{figure}

 These results demonstrate good agreement between the solutions obtained by different methods, with the exception of the frequency at which resonance occurs. 
 
 The influence of the angle \(\beta\) on the frequency dependence of the SIF is illustrated in Figs.~\ref{PopovFig12} and \ref{PopovFig13}. Figure~\ref{PopovFig12} shows the frequency dependence of the SIF for normal stresses under normal crack loading, while Fig.~\ref{PopovFig13} presents the corresponding dependence for tangential stresses under shear loading, for various values of the angle \(\beta\).
 
 In the case of normal loading, all angle values \(\beta\) exhibit a global maximum in the SIF, with the frequency at which this maximum occurs depending on the angle. The highest peak value of the SIF is observed at \(\beta = 0^\circ\); however, at high frequencies \(\kappa_0 > 1.8\), the largest SIF values occur for \(\beta = 30^\circ\). Under shear loading, the SIF for tangential stresses tends to decrease with increasing frequency and shows only minor deviations from the corresponding SIF of a single crack.
 
 \begin{figure}[b]
 	\sidecaption[b]
 	\includegraphics[width=.45\textwidth]{31_Popov/PopovFig12.jpg}
 	\caption{SIF of normal stresses under normal loading of cracks.}
 	\label{PopovFig12}
 \end{figure}
 
 \begin{figure}[b]
 	\sidecaption[b]
 	\includegraphics[width=.45\textwidth]{31_Popov/PopovFig13.jpg}
 	\caption{SIF of tangential stresses under shear loading of cracks.}
 	\label{PopovFig13}
 \end{figure}
 
 To investigate in greater detail the influence of crack interaction on the SIF under dynamic loading, the system of nine cracks depicted in Fig.~\ref{PopovFig14} was analyzed.
 
 \begin{figure}[t]
 	\sidecaption[t]
 	\includegraphics[width=.45\textwidth]{31_Popov/PopovFig14.png}
 	\caption{System of nine cracks.}
 	\label{PopovFig14}
 \end{figure}
 
 The configuration of the crack system is defined by the parameters $\theta$, $h$, and $\beta$. In the present calculations, the values $\theta = h = d_1/3$ and $d_k = d_1/3$ ($k = 2, \dots, 9$) were used. The system was subjected to normal loading, and the results are presented in Fig.~\ref{PopovFig15} as graphs illustrating the frequency dependence of the SIF of normal stresses, $|k_1^{+}|$, for the first crack at various orientation angles $\beta$.
 
 \begin{figure}[t]
 	\sidecaption[t]
 	\includegraphics[width=.45\textwidth]{31_Popov/PopovFig15.jpg}
 	\caption{Influence of the crack orientation angle $\beta$ on the SIF.}
 	\label{PopovFig15}
 \end{figure}
 
 As evident from Fig.~\ref{PopovFig15}, the configuration of the crack system significantly affects the behavior of the SIF in the frequency domain. A global maximum is observed in each case, with the highest value occurring at $\beta = 0^\circ$, corresponding to the configuration where all cracks are parallel. Under static loading and at low frequencies ($\kappa_0 < 1.2$), the presence of neighboring cracks reduces the SIF near the crack tip.
 
 To further examine the effect of crack interaction, Fig.~\ref{PopovFig16} shows the influence of the number of interacting cracks on the SIF of the first crack. The following configurations were considered: two cracks (first and second), three (first through third), five, seven, and all nine cracks (see Fig.~\ref{PopovFig14}).
 
 \begin{figure}[b]
 	\sidecaption[b]
 	\includegraphics[width=.45\textwidth]{31_Popov/PopovFig16.jpg}
 	\caption{Influence of the number of interacting cracks on the SIF.}
 	\label{PopovFig16}
 \end{figure}
 
 According to Fig.~\ref{PopovFig16}, increasing the number of cracks from one to five leads to a marked increase in the peak SIF values. However, further increases beyond five cracks produce only marginal changes in the maximum SIF. In all considered cases, the presence of adjacent cracks reduces the SIF at low frequencies ($\kappa_0 < 1.2$).
  	
	
\section{Stress State and Stress Intensity Factors in a Cylindrical Body with Cracks under Dynamic Surface Loading}

Consider a self-balanced harmonic normal load of the form $\mu P\left(\varphi \right)e^{-i\omega t}$ applied to the lateral surface of an infinite cylindrical body with cracks. The cross-section of the cylinder is depicted in Fig.~\ref{PopovFig06}. Under these loading conditions, the cylinder experiences a plane strain state. Accordingly, only the displacement components $u_{r}$ and $u_{\varphi}$ need to be determined, satisfying the Lamé equations:

\begin{equation} \label{PopovEq857539}
	\begin{aligned}
		\mu \left(\Delta u_{r} -\dfrac{u_{r}}{r^{2}} - \dfrac{2}{r^{2}} \dfrac{\partial u_{\varphi}}{\partial \varphi} \right)
		+ (\lambda + \mu) \dfrac{\partial}{\partial r} \left(\dfrac{1}{r} \dfrac{\partial}{\partial r}(r u_{r}) + \dfrac{1}{r} \dfrac{\partial u_{\varphi}}{\partial \varphi} \right)
		+ \rho \omega^{2} u_{r} &= 0;
		\\
		\mu \left(\Delta u_{\varphi} -\dfrac{u_{\varphi}}{r^{2}} + \dfrac{2}{r^{2}} \dfrac{\partial u_{r}}{\partial \varphi} \right)
		+ (\lambda + \mu) \dfrac{\partial}{\partial \varphi} \left(\dfrac{1}{r} \dfrac{\partial}{\partial r}(r u_{r}) + \dfrac{1}{r} \dfrac{\partial u_{\varphi}}{\partial \varphi} \right)
		+ \rho \omega^{2} u_{\varphi} &= 0.
	\end{aligned}
\end{equation}

The corresponding boundary conditions on the lateral surface of the cylinder are:

\begin{equation} \label{PopovEq430273}
	\begin{array}{rcl}
		\sigma_{xx} \cos(\bar{n}, x) + \sigma_{xy} \cos(\bar{n}, y) &=& \mu P(\varphi) \cos(\bar{n}, x), \\
		\sigma_{xy} \cos(\bar{n}, x) + \sigma_{yy} \cos(\bar{n}, y) &=& \mu P(\varphi) \cos(\bar{n}, y),
	\end{array}
\end{equation}
where $(x, y) \in L$ (with $L$ being the cross-sectional contour), $\bar{n}$ is the outward unit normal vector to the surface, and $\sigma_{xx}$, $\sigma_{yy}$, and $\sigma_{xy}$ denote the Cartesian components of the stress tensor. These are related to the polar components $\sigma_{rr}$, $\sigma_{\varphi\varphi}$, and $\sigma_{r\varphi}$ through the following transformation:

\begin{equation} \label{PopovEq616969}
	\left(
	\begin{array}{c}
		\sigma_{yy} \\
		\sigma_{xx} \\
		\sigma_{xy}
	\end{array}
	\right)
	=
	\left(
	\begin{array}{ccc}
		\sin^{2} \varphi & \cos^{2} \varphi & 2 \sin \varphi \cos \varphi \\
		\cos^{2} \varphi & \sin^{2} \varphi & -2 \sin \varphi \cos \varphi \\
		\sin \varphi \cos \varphi & -\sin \varphi \cos \varphi & \cos^{2} \varphi - \sin^{2} \varphi
	\end{array}
	\right)
	\left(
	\begin{array}{c}
		\sigma_{rr} \\
		\sigma_{\varphi\varphi} \\
		\sigma_{r\varphi}
	\end{array}
	\right).
\end{equation}

To formulate the boundary conditions on the crack surfaces, a local coordinate system $O_k x_k y_k$ is associated with each crack (see Fig.~\ref{PopovFig06}). Let $u_k(x_k, y_k)$, $v_k(x_k, y_k)$, $\tau_{kxy}(x_k, y_k)$, $\sigma_{kx}(x_k, y_k)$, and $\sigma_{ky}(x_k, y_k)$ denote the displacement and stress components in the local coordinate system associated with the $k$th crack. The crack surfaces are traction-free, which implies the following boundary conditions:

\begin{equation} \label{PopovEq953877}
	\sigma_{ky}(x_k, 0) = 0,\quad \tau_{kxy}(x_k, 0) = 0,\quad -d_k < x_k < d_k\quad (k=1,2,\dots,N).
\end{equation}

Furthermore, displacement discontinuities with prescribed jumps are assumed along the crack surfaces, as described by \eqref{PopovEq200550}. The local stress components in the $O_k x_k y_k$ system are related to the polar stress components via the transformation:

\begin{equation} \label{PopovEq604994}
	\left(\begin{array}{c}
		\sigma_{ky} \\
		\sigma_{kx} \\
		\sigma_{kxy}
	\end{array}\right)
	=
	\left(\begin{array}{ccc}
		\sin^2(\varphi - \alpha_k) & \cos^2(\varphi - \alpha_k) & \sin 2(\varphi - \alpha_k) \\
		\cos^2(\varphi - \alpha_k) & \sin^2(\varphi - \alpha_k) & -\sin 2(\varphi - \alpha_k) \\
		\frac{1}{2} \sin 2(\varphi - \alpha_k) & -\frac{1}{2} \sin 2(\varphi - \alpha_k) & \cos 2(\varphi - \alpha_k)
	\end{array}\right)
	\left(\begin{array}{c}
		\sigma_{rr} \\
		\sigma_{\varphi\varphi} \\
		\sigma_{r\varphi}
	\end{array}\right).
\end{equation}

Thus, the steady-state dynamic problem of a cracked body is reduced to solving the governing equations \eqref{PopovEq857539} with boundary conditions \eqref{PopovEq430273}, \eqref{PopovEq953877}, and \eqref{PopovEq200550}. For this purpose, discontinuous solutions of the equations of motion in local coordinates—such as \eqref{PopovEq144429} and \eqref{PopovEq460990}, featuring the prescribed displacement jumps—are employed.

The total displacement field in the polar coordinate system is represented as:

\begin{equation} \label{PopovEq481837}
	\begin{aligned}
		u_r(r, \varphi) &= u_{0r}(r, \varphi) + \displaystyle\sum_{k=1}^{N} u_{dkr}(r, \varphi), \\
		u_{\varphi}(r, \varphi) &= u_{0\varphi}(r, \varphi) + \displaystyle\sum_{k=1}^{N} u_{dk\varphi}(r, \varphi),
	\end{aligned}
\end{equation}
where $u_{dkr}$ and $u_{dk\varphi}$ are obtained by transforming the discontinuous solutions from the local to the global (polar) system using:

\begin{equation} \label{31.46)}
	\left(\begin{array}{c}
		u_{dlr} \\
		u_{dl\varphi}
	\end{array}\right)
	=
	\left(\begin{array}{cc}
		\cos(\varphi - \alpha_l) & \sin(\varphi - \alpha_l) \\
		-\sin(\varphi - \alpha_l) & \cos(\varphi - \alpha_l)
	\end{array}\right)
	\left(\begin{array}{c}
		u_{dl} \\
		v_{dl}
	\end{array}\right).
\end{equation}

The functions $u_{0r}(r, \varphi)$ and $u_{0\varphi}(r, \varphi)$ represent the unknown smooth solution to Eq.~\eqref{PopovEq857539}, required to satisfy the surface boundary condition \eqref{PopovEq430273}. They are expressed via wave potentials \cite{Popov1} as:

\begin{equation} \label{31.47)}
	u_{0r} = \dfrac{\partial \psi_1}{\partial r} + \dfrac{1}{r} \dfrac{\partial \psi_2}{\partial \varphi}, \qquad
	u_{0\varphi} = \dfrac{1}{r} \dfrac{\partial \psi_1}{\partial \varphi} - \dfrac{\partial \psi_2}{\partial r},
\end{equation}
where $\psi_1$ and $\psi_2$ satisfy the Helmholtz equations:

\begin{equation} \label{PopovEq947656}
	\Delta \psi_j + \kappa_j^2 \psi_j = 0, \quad \kappa_j^2 = \dfrac{\omega^2}{c_j^2}, \;\;
	(j = 1,2), \qquad
	c_1^2 = \dfrac{\lambda + 2\mu}{\rho}, \;\; c_2^2 = \dfrac{\mu}{\rho}.
\end{equation}

The corresponding stress components in polar coordinates are then given by:

\begin{equation} \label{PopovEq359974}
	\begin{aligned}
		\sigma_{0rr} &= -\mu \left( \kappa_2^2 \psi_1 + \dfrac{2}{r^2} \dfrac{\partial^2 \psi_1}{\partial \varphi^2}
		+ \dfrac{2}{r} \dfrac{\partial \psi_1}{\partial r}
		+ \dfrac{2}{r^2} \dfrac{\partial \psi_2}{\partial \varphi}
		- \dfrac{2}{r} \dfrac{\partial^2 \psi_2}{\partial r \partial \varphi} \right),
		\\
		\sigma_{0\varphi\varphi} &= -\mu \left( \kappa_2^2 \psi_1
		+ 2 \dfrac{\partial^2 \psi_1}{\partial r^2}
		- \dfrac{2}{r^2} \dfrac{\partial \psi_2}{\partial \varphi}
		+ \dfrac{2}{r} \dfrac{\partial^2 \psi_2}{\partial r \partial \varphi} \right),
		\\
		\sigma_{0r\varphi} &= \mu \left(
		\dfrac{2}{r} \dfrac{\partial^2 \psi_1}{\partial r \partial \varphi}
		- \dfrac{2}{r^2} \dfrac{\partial \psi_1}{\partial \varphi}
		+ \kappa_2^2 \psi_2
		+ \dfrac{2}{r^2} \dfrac{\partial^2 \psi_2}{\partial \varphi^2}
		+ \dfrac{2}{r} \dfrac{\partial \psi_2}{\partial r}
		\right).
	\end{aligned}
\end{equation}

	The wave potentials are sought in the form of linear combinations of $M$ partial solutions of the Helmholtz equations \eqref{PopovEq947656}, analogous to \eqref{PopovEq659953}:
	
	\begin{equation} \label{PopovEq299502}
		\psi_1 = r_0^2 \sum_{i=1}^{M} A_i\, g_{1i}(r, \varphi), \quad
		\psi_2 = r_0^2 \sum_{i=1}^{M} B_i\, g_{2i}(r, \varphi),
	\end{equation}
	where the basis functions $g_{si}(r, \varphi)$ are defined as:
	\[
	\begin{array}{rcl}
		g_{s(2m-1)}(r,\varphi) &=& J_{m-1}(\kappa_i r) \cos[(m-1)\varphi], \\
		g_{s(2m)}(r,\varphi) &=& J_m(\kappa_i r) \sin(m\varphi), \qquad (s = 1, 2),
	\end{array}
	\]
	with $J_m$ denoting the Bessel functions of the first kind.
	
	Substituting these expressions into the stress formulas \eqref{PopovEq359974} yields the following expressions for the stress components:
	
	\begin{equation} \label{PopovEq955753} 
		\begin{aligned} \sigma _{0rr} &=
			\begin{multlined}[t]
				-\mu r_{0}^{2} \Bigg(\displaystyle\sum _{i=1}^{M}A_{i} \left(\kappa _{2}^{2} g_{1i} \left(r,\varphi \right)+\frac{2}{r^{2} } \frac{\partial ^{2} g_{1i} \left(r,\varphi \right)}{\partial \varphi ^{2} } +\frac{2}{r} \frac{g_{2i} \left(r,\varphi \right)}{\partial r} \right) 
				\\
				+\displaystyle\sum _{i=1}^{M}B_{i} \left(\frac{2}{r^{2} } \frac{\partial g_{1i} \left(r,\varphi \right)}{\partial \varphi } -\frac{2}{r} \frac{\partial ^{2} g_{2i} \left(r,\varphi \right)}{\partial r\partial \varphi } \right) \Bigg);
			\end{multlined}
			\\
			\sigma _{0\varphi \varphi } &=
			\begin{multlined}[t]
				-\mu r_{0}^{2} \Bigg(\displaystyle\sum _{i=1}^{M}A_{i} \left(\kappa _{2}^{2} g_{1i} \left(r,\varphi \right)+2\frac{\partial ^{2} g_{1i} \left(r,\varphi \right)}{\partial r^{2} } \right) 
				\\
				+\displaystyle\sum _{i=1}^{M}B_{i} \left(\frac{2}{r} \frac{\partial ^{2} g_{2i} \left(r,\varphi \right)}{\partial r\partial \varphi } -\frac{2}{r^{2} } \frac{\partial g_{2i} \left(r,\varphi \right)}{\partial \varphi } \right) \Bigg); 
			\end{multlined}
			\\
			\sigma _{0r\varphi } &=
			\begin{multlined}[t]
				\mu r_{0}^{2} \Bigg(\displaystyle\sum _{i=1}^{M}A_{i} \left(\frac{2}{r} \frac{\partial ^{2} g_{1i} \left(r,\varphi \right)}{\partial r\partial \varphi } -\frac{2}{r^{2} } \frac{\partial g_{1i} \left(r,\varphi \right)}{\partial \varphi } \right)
				\\
				+\displaystyle\sum _{i=1}^{M}B_{i}  \left(\kappa _{2}^{2} g_{2i} \left(r,\varphi \right)+\frac{2}{r^{2} } \frac{\partial ^{2} g_{2i} \left(r,\varphi \right)}{\partial \varphi ^{2} } +\frac{2}{r} \frac{\partial g_{2i} \left(r,\varphi \right)}{\partial r} \right)\Bigg).  
			\end{multlined}
		\end{aligned}
	\end{equation} 
	
Next, the unknown displacement jumps \eqref{PopovEq200550} are to be determined from the boundary conditions \eqref{PopovEq953877} on the crack surfaces, and the coefficients $A_i$ and $B_i$ are to be found from the boundary conditions \eqref{PopovEq430273} on the cylinder surface. 

To apply the boundary conditions \eqref{PopovEq953877}, the stresses in equation \eqref{PopovEq491412} are expressed using the representation \eqref{PopovEq481837} as follows:

\begin{equation} \label{PopovEq415691}
	\begin{aligned}
		\sigma_{ky}(x_k, y_k) &= \sigma_{0y}(x_k, y_k) + \displaystyle\sum\limits_{l=1}^{N} \sigma_{kl}(x_k, y_k), \\
		\tau_{kyx}(x_k, y_k) &= \tau_{0yx}(x_k, y_k) + \displaystyle\sum\limits_{l=1}^{N} \tau_{kl}(x_k, y_k),
	\end{aligned}
\end{equation}
where $\sigma_{kl}(x_k, y_k)$ and $\tau_{kl}(x_k, y_k)$ are derived from the stresses caused by displacement jumps \eqref{PopovEq460990} and are transformed into local coordinates using \eqref{PopovEq491412}:

\[
\begin{array}{rcl}
	\sigma_{kl} &=& \sigma_{dx_l} \sin^2 \alpha_{kl} + \sigma_{dy_l} \cos^2 \alpha_{kl} - \tau_{dy_l x_l} \sin 2\alpha_{kl}, \\
	2\tau_{kl} &=& -\sigma_{dx_l} \sin 2\alpha_{kl} + \sigma_{dy_l} \sin 2\alpha_{kl} + 2\tau_{dy_l x_l} \cos 2\alpha_{kl}, \quad \alpha_{kl} = \alpha_k - \alpha_l.
\end{array}
\]

The stresses $\sigma_{0y}(x_k, y_k)$ and $\tau_{0yx}(x_k, y_k)$ are derived from the expressions \eqref{PopovEq955753} using the stress transformation formulas \eqref{PopovEq604994}. Substituting the expressions \eqref{PopovEq415691} into the boundary conditions \eqref{PopovEq953877} yields a system of integro-differential equations for the displacement jumps:

\begin{equation} \label{PopovEq949735}
	\varphi_{sk}(\tau) = d_k^{-1} \chi_{sk}(d_k \tau) \quad (s = 3, 4; \; k = 1, 2, \dots, N),
\end{equation}
analogous to system \eqref{PopovEq123113}, with the right-hand sides defined by

\begin{equation} \label{PopovEq882335}
	\begin{aligned}
		f_k(\zeta) &= \displaystyle\sum_{i=1}^{M} \left( A_i F_{ik}(\zeta) + B_i F_{i+M,k}(\zeta) \right), \\
		h_k(\zeta) &= \displaystyle\sum_{i=1}^{M} \left( A_i H_{ik}(\zeta) + B_i H_{i+M,k}(\zeta) \right).
	\end{aligned}
\end{equation}

\noindent The functions $F_{ik}(\zeta)$ and $H_{ik}(\zeta)$ $(i = 1, \dots, 2M;\; k = 1, \dots, N)$ are expressed through the functions $g_{1i}(r, \varphi)$, $g_{2i}(r, \varphi)$ and their derivatives.

Due to the linearity of the system and the structure of the right-hand sides \eqref{PopovEq882335}, the solution of \eqref{PopovEq949735} can be sought in the form:

\begin{equation} \label{31.55)}
	\varphi_{sk}(\tau) = \sum_{i=1}^{M} \left( A_i \varphi_{isk}(\tau) + B_i \varphi_{i+M,sk}(\tau) \right) \quad (s = 3, 4; \; k = 1, 2, \dots, N).
\end{equation}

\noindent Then, the functions $\varphi_{isk}(\tau)$ are found from the following systems of integro-differential equations:

\begin{subequations} \label{PopovEq639348}
\begin{multline} \label{PopovEq639348a} 
			\frac{1}{2\pi } \int _{-1}^{1}\varphi '_{i3k} (\tau ) \left[-\frac{2\left(1-\xi ^{2} \right)}{\tau -\xi } +F_{13kk} \left(\tau -\xi \right)\right]d\tau 
			\\
			+\frac{1}{2\pi } \int _{-1}^{1}\varphi _{i3k} (\tau ) \left[\kappa _{0}^{2} \gamma _{k}^{2} \ln \left|\tau -\varsigma \right|+U_{13kk} \left(\tau -\varsigma \right)\right]d\tau 
			\\
			+\frac{1}{2\pi }\sum_{\substack{l=1 \\ l \ne k}}^{N}\Bigg[ \int _{-1}^{1}\varphi '_{i3l} (\tau )F_{13kl} (\tau ,\varsigma )d\tau 
			+ \int _{-1}^{1}\varphi '_{i4k} (\tau )F_{14kl} (\tau ,\varsigma ) d\tau  \Bigg] 
			\\
			+\frac{1}{2\pi }\sum_{\substack{l=1 \\ l \ne k}}^{N}\Bigg[ \int _{-1}^{1}\varphi _{i3l} (\tau )U_{13kl} (\tau ,\varsigma )d\tau 
			+ \int _{-1}^{1}\varphi _{i4k} (\tau )U_{14kl} (\tau ,\varsigma ) d\tau  \Bigg] =F_{ik} (\zeta );
		\end{multline} 
		
		\begin{multline}\label{PopovEq639348b} 
			\frac{1}{2\pi } \int _{-1}^{1}\varphi '_{i4k} (\tau ) \left[\frac{2\left(1-\xi ^{2} \right)}{\tau -\varsigma } +F_{24kk} \left(\tau -\varsigma \right)\right]d\tau 
			\\
			+\frac{1}{2\pi } \int _{-1}^{1}\varphi _{i4k} (\tau ) \left[\kappa _{0}^{2} \gamma _{k}^{2} \ln \left|\tau -\varsigma \right|+U_{24kk} \left(\tau -\varsigma \right)\right]d\tau 
			\\
			+\frac{1}{2\pi }\sum_{\substack{l=1 \\ l \ne k}}^{N}\Bigg[ \int _{-1}^{1}\varphi '_{i3l} (\tau )F_{23kl} (\tau ,\varsigma )d\tau 
			+\int _{-1}^{1}\varphi '_{i4k} (\tau )F_{24kl} (\tau ,\varsigma ) d\tau  \Bigg] 
			\\
			+\frac{1}{2\pi }\sum_{\substack{l=1 \\ l \ne k}}^{N}\Bigg[ \int _{-1}^{1}\varphi _{i3l} (\tau )U_{23kl} (\tau ,\varsigma )d\tau 
			+ \int _{-1}^{1}\varphi _{i4k} (\tau )U_{24kl} (\tau ,\varsigma ) d\tau  \Bigg] =H_{ik} (\zeta )
		\end{multline}
	\end{subequations}
where $k = 1, 2, \dots, N$ and $i = 1, 2, \dots, 2M$.

The numerical solution of system \eqref{PopovEq639348} is carried out analogously to the solution of system \eqref{PopovEq143660}. The derivatives of the unknown functions $\varphi_{isk}(\tau)$, accounting for their singular behavior, are approximated using representations \eqref{PopovEq992786} and \eqref{PopovEq972705}.

To compute $\psi_{iskm} = \psi_{isk}(\tau_m)$, a system of linear algebraic equations is derived from \eqref{PopovEq639348} by applying the method of mechanical quadratures.

Once system \eqref{PopovEq639348} is solved, the remaining task is to determine the unknown coefficients $A_i$ and $B_i$. To this end, the boundary values of the stresses are computed as:

\begin{equation} \label{PopovEq185218}
	\sigma_{xx} = \sigma_{0xx} + \sum_{l=1}^{N} \sigma_{lxx}, \quad
	\sigma_{yy} = \sigma_{0yy} + \sum_{l=1}^{N} \sigma_{lyy}, \quad
	\tau_{yx} = \tau_{0yx} + \sum_{l=1}^{N} \tau_{lyx},
\end{equation}
where the components $\sigma_{lxx}$, $\sigma_{lyy}$, and $\tau_{lyx}$ are derived from the stresses of the discontinuous solutions \eqref{PopovEq460990} using formulas \eqref{PopovEq642436}, and the terms $\sigma_{0xx}$, $\sigma_{0yy}$, and $\tau_{0yx}$ are obtained by substituting the expressions \eqref{PopovEq955753} into the transformation relations \eqref{PopovEq616969}.

By applying the collocation method at the angular points $\varphi_p = 2\pi(p - 1)/M$, $p = 1, 2, \dots, M$, the boundary conditions \eqref{PopovEq430273} reduce to a system of $2M$ linear equations for determining the coefficients $A_i$ and $B_i$.

After solving the integral equation system \eqref{PopovEq143660} and obtaining the coefficients, the approximate values of the stress intensity factors (SIFs) for each crack are calculated using the formulas:

\begin{equation} \label{31.58)}
	K_{(s-2)l}^{\pm} = \mu \sqrt{d_l} \, k_{(s-2)l}^{\pm} \quad (l = 1, 2, \dots, N;\; s = 3, 4),
\end{equation}

\begin{align*}
	k_{(s-2)l}^{+} &= -\dfrac{1 - \xi^2}{n} \displaystyle\sum_{m=1}^{n} (-1)^m \psi_{slm} \cot\left( \dfrac{\beta_m}{2} \right), \\[1ex]
	k_{(s-2)l}^{-} &= -\dfrac{1 - \xi^2}{n} (-1)^{s+n} \displaystyle\sum_{m=1}^{n} (-1)^m \psi_{slm} \tan\left( \dfrac{\beta_m}{2} \right),
\end{align*}

\[
\psi_{slm} = \sum_{i=1}^{M} \left( A_i \psi_{islm} + B_i \psi_{i+M,slm} \right).
\]

To demonstrate the numerical implementation of the proposed method, we consider an elliptical cylinder with eccentricity $\varepsilon = 0.5$ subjected to a self-balanced normal load $P(\varphi) = 1$, and containing a single crack with its tip located at the origin (Fig.~\ref{PopovFig17}).

\begin{figure}[!h]
	\sidecaption[b]
	\includegraphics[width=.45\textwidth]{31_Popov/PopovFig17.png}
	\caption{Cross-section of the elliptical cylinder with a crack.}
	\label{PopovFig17}
\end{figure}

Figures~\ref{PopovFig18} and~\ref{PopovFig19} present the frequency response of the absolute values of the normal and shear SIFs for various crack inclination angles, when the crack tip is located near the boundary. Curves 1, 2, 3, and 4 correspond to crack orientations of $\alpha = 0^\circ$, $45^\circ$, $60^\circ$, and $90^\circ$, respectively, with a fixed relative crack length $\gamma = d/r_0 = 0.4$.

\begin{figure}[t]
	\sidecaption[t]
	\includegraphics[width=.5\textwidth]{31_Popov/PopovFig18.png}
	\caption{SIF of normal (breakaway) stresses for different crack angles.}
	\label{PopovFig18}
\end{figure}

\begin{figure}[t]
	\sidecaption[t]
	\includegraphics[width=.5\textwidth]{31_Popov/PopovFig19.png}
	\caption{SIF of shear stresses for different crack angles.}
	\label{PopovFig19}
\end{figure}

A notable feature in the SIF behavior is the occurrence of local maxima at frequencies below the first resonance. The frequency and magnitude of these peaks are strongly influenced by the crack inclination angle.

	
	\section{Conclusions}
	
	An iterative method has been proposed for determining the stress intensity factors (SIFs) in systems of dynamically loaded interacting cracks. This approach eliminates the need to solve large-scale systems of integro-differential equations. At each iteration step, a set of independent equations corresponding to individual cracks is solved. Numerical studies have demonstrated good agreement between the iterative method and the direct solution of the full system of integro-differential equations. The method exhibits both stability and convergence, even for system of densely spaced cracks of complex configuration.
	
	A key result from the fracture mechanics perspective is the identification of absolute maxima in the SIF values at specific frequencies of oscillatory loading. These peak values can significantly exceed (by several times) the corresponding SIFs for isolated cracks under similar loading conditions. Conversely, under static or low-frequency loading, the SIF values may be reduced compared to those for individual cracks due to shielding effects.
	
	Furthermore, a method has been developed to solve two-dimensional elasticity problems in the frequency domain for finite regions bounded by arbitrarily smooth contours and containing cracks. This method relies on representing the displacements and stresses in the frequency domain using discontinuous solutions to the equations of motion for elastic cylindrical bodies of arbitrary cross-section under longitudinal shear or plane strain conditions. An additional function is introduced to account for the oscillatory influence of the external surface. As a result, the boundary conditions on the cracks and on the body's surface are satisfied sequentially and independently.
	
	The resulting systems of integro-differential equations derived from the conditions on the cracks are independent of the body's boundary shape and differ only in their right-hand sides. Numerical analysis of the SIF behavior near tunnel cracks in cylindrical bodies as a function of the dimensionless wave number reveals several important trends. The frequency response of the SIFs is highly sensitive to the crack location, size, and proximity to the boundary. It has also been shown that the presence of cracks affects the resonant frequencies of the body: for a fixed geometry, the first resonant frequency decreases as the crack approaches the loaded boundary.
	
	The proposed methods are generalizable and can be extended to systems with other types of defects, particularly thin rigid inclusions.
	
	
\bibliographystyle{unsrtnat}
\bibliography{Chapter31}
	
	
	
